\makeabstract
    {
    You Spin Me Right {\textquoteleft}Round, Baby, Right {\textquoteleft}Round: Accounting for Non-Inertial Effects in the Search for Long-Range Spin-Dependent Interactions
    }
    {
    Daniel Murphy\textsuperscript{1}, August Castelli\textsuperscript{1}, Nathan Clayburn\textsuperscript{1,2}, Andrew Glassford\textsuperscript{1}, Ashley Kim\textsuperscript{1}, Larry Hunter\textsuperscript{1}
    }
    {
    \textsuperscript{1} Department of Physics \& Astronomy, Amherst College, Amherst, MA\\
\textsuperscript{2} Department of Physics, Keene State College, Keene, NH
    }
    {
    We search for Long-Range Spin-Dependent Interactions (LRSDIs), which are predicted by some Standard Model extensions and serve to limit potential dark matter candidates. Within a homogeneous magnetic field, we spin-polarize Hg-199 and Cs-133 and measure their Larmor precession frequencies {\textendash} the difference is our signal. We vary the orientation of our apparatus, and thus magnetic field, and look for changes in this signal. Because LRSDIs are orientation dependent, such changes could be a signature of one of these interactions. However, as a result of Earth's rotation, we are making these measurements in a non-inertial reference frame, which causes an orientation dependent frequency shift in our signal. To properly account for this frequency shift, we need to ensure that our applied magnetic field is closely aligned to true cardinal directions in every table position and measure its misalignment.
    }
    {
    Our lab has a true North reference with a precision of 0.042 degrees, to which we attach strings. We photograph the strings and magnetic field coils within the apparatus, then manually feed their positions into a Python program which outputs the angle difference between the strings and coils {\textendash} our misalignment. We then adjust a positioning mirror on the apparatus accordingly, realign using our laser positioning system, and repeat. We aim for a misalignment of zero degrees within a 95\% confidence interval, accounting for the uncertainty of the North reference.
    }
    {
    After performing the procedure described, we achieved North and South misalignments of -0.014{\textdegree} \ensuremath{\pm} 0.025{\textdegree}  and -0.003{\textdegree} \ensuremath{\pm} 0.022{\textdegree} respectively and East and West misalignments of 0.013{\textdegree} \ensuremath{\pm} 0.022{\textdegree} and 0.011{\textdegree} \ensuremath{\pm} 0.016{\textdegree} respectively. We account for the measured misalignment from each cardinal direction in our final data by subtracting the associated frequency shift calculated by finding the component of Earth{\textquoteright}s spin axis which is in the direction of the applied magnetic. Since East-West frequency shifts are dependent on the sine of the angle between our applied field and true East, we are first order sensitive to misalignment for these directions. We are only second order sensitive to misalignment for North-South, which are cosine dependent. The uncertainty in our misalignment measurements also translates to a nHz Hg uncertainty which is added in quadrature to our other systematic uncertainties.
    }
    {
    Having properly aligned our apparatus, we are now able to account for misalignment in all cardinal directions and have paved the way for future data collection.
    }

    \newpage
    